% §6.1 Shen-Ross uniqueness channel, clean 12-city panel.
% Numbers locked against results/replications/shen_12city_table.csv,
% shen_12city_pooled.json, and type_s_type_m_shen.csv. Supersedes v2, whose
% per-city thetas, "two markets significant", n approx 285-334, 25% power, and
% Type-M factor of 2 were all stale, and whose Portland-degenerate / k=11-vs-k=12
% framing no longer applies. Flowing prose, no em dashes, no enumerations.

\subsection{Shen-Ross uniqueness channel across twelve metropolitan markets}
\label{sec:shen-12city}

The Shen-Ross uniqueness channel \citep{shen2021information} measures how far the
language in a listing departs from the language its neighbors use, quantifying
departure as the cosine distance between each listing's vector-space text
representation and the centroid of its $K = 5$ nearest peers in
latitude-longitude space. We apply the resulting per-listing uniqueness score as
the treatment in a partially-linear DML estimator with cross-fitted ridge
nuisances solved through a Cholesky factorization, following the Shen-Ross
uniqueness construction as a close adaptation rather than a verbatim replication
of the original Doc2Vec and MLS-area-by-year cell pipeline. Effects are reported
per standard deviation of the uniqueness score, and the $95\%$ confidence
intervals are formed from the cross-fitting influence-function standard error.

Table~\ref{tab:shen-12city} reports the per-$\sigma$ uniqueness effect for each
of the twelve metropolitan markets, together with the influence-function
confidence interval and the indicator for whether that interval excludes zero.
The panel is now clean throughout, with no degenerate market and no exclusion,
and the per-market samples range from $986$ deduplicated listings in San
Francisco to $15{,}227$ in New York, an order of magnitude larger than the
corpus on which the earlier draft was built. The DML interval excludes zero in
ten of the twelve markets. Eight of those ten identify a positive uniqueness
premium. The effect is largest in Philadelphia at $+0.191$ with interval
$[+0.168, +0.213]$ and in San Francisco at $+0.111$ with interval
$[+0.063, +0.159]$, followed by Chicago at $+0.074$ $[+0.044, +0.103]$, and then
a tight cluster of mid-sized markets, with Denver at $+0.045$
$[+0.027, +0.062]$, Phoenix at $+0.041$ $[+0.021, +0.060]$, Dallas at $+0.041$
$[+0.023, +0.059]$, Atlanta at $+0.037$ $[+0.018, +0.056]$, and Portland at
$+0.029$ $[+0.008, +0.050]$. Atlanta itself, the market on which Shen's
identification was originally established, returns a positive and significant
per-$\sigma$ effect of $+0.037$, roughly a quarter of the magnitude of her
published Atlanta estimate but of the same sign and recovered under a far more
conservative confounder set. The remaining two of the ten significant markets
carry the opposite sign, with New York at $-0.024$ $[-0.036, -0.013]$ and
Washington at $-0.031$ $[-0.054, -0.007]$, and we record these sign reversals
rather than absorb them into the positive consensus. The final two markets,
Boston at $-0.004$ $[-0.031, +0.024]$ and Seattle at $+0.016$
$[-0.011, +0.043]$, return intervals that straddle zero.

The twelve per-market estimates pool to a positive effect. A random-effects
combination places the pooled uniqueness premium at $+0.043$ per standard
deviation with $95\%$ confidence interval $[+0.010, +0.076]$, and the
precision-weighted fixed-effect combination sits at $+0.029$ with a standard
error of $0.003$. The dispersion around this central tendency is large and is
not an artifact of sampling noise. The heterogeneity statistic is $I^2 = 96.7\%$
with Cochran's $Q = 335.9$ on eleven degrees of freedom and a between-market
variance of $\tau^2 = 0.0033$, so almost all of the variation across markets
reflects genuine between-market differences in the operative effect rather than
estimation error within markets. Section~\ref{sec:meta-regression-12} takes up
this heterogeneity directly with a random-effects meta-regression using
Paule-Mandel $\tau^2$ and Hartung-Knapp-Sidik-Jonkman small-sample intervals,
and we defer the moderator analysis to that section rather than adjudicate it
here.

\paragraph{Design analysis.} The cross-market pattern cannot be read as an
artifact of low power, because the clean panel is well powered. A Gelman-Carlin
\citep{gelman2014beyond} retrodesign of each per-market estimator against an
assumed true effect of $+0.050$ per standard deviation, a benchmark set
deliberately below Shen's published Atlanta estimate, returns a median
per-market power of $0.99$ across the twelve markets, far above the conventional
$0.80$ target. The Type-S sign-error rate is negligible everywhere, with the
largest value across the panel below $6 \times 10^{-5}$, so the sign of any
detected effect is reliable. The median Type-M exaggeration factor is $1.00$,
indicating essentially no magnitude inflation in the typical market. The single
exception is San Francisco, the smallest sample at $986$ listings and the only
market with retrodesign power below $0.80$, where power is $0.54$ and the Type-M
factor is $1.36$. The San Francisco point estimate should therefore be read as
mildly exaggerated in magnitude, and its influence-function interval, resting on
the smallest sample in the panel, carries the mild anti-conservatism that the
analytic variance can show near a thousand observations, a feature we document
in the coverage simulation. Everywhere else the design analysis supports reading
both the significant detections and the heterogeneity across markets at face
value.

\paragraph{Magnitudes.} Because the outcome is log sale price, the per-$\sigma$
coefficients translate directly into percentage differences in price. A
one-standard-deviation increase in listing uniqueness corresponds to a price
difference of roughly $+21.0\%$ in Philadelphia, $+11.7\%$ in San Francisco,
$+7.7\%$ in Chicago, and a cluster near $+4\%$ across Denver, Dallas, Phoenix,
and Atlanta, falling to $+3.0\%$ in Portland, with the two significant reversals
at $-2.4\%$ in New York and $-3.0\%$ in Washington. The external anchor for these
figures is the published Shen-Ross Atlanta estimate of $+0.149$ per standard
deviation with a standard error of $0.034$, obtained on a sample of roughly forty
thousand listings under MLS-area-by-year fixed effects. Our Atlanta estimate of
$+0.037$ recovers the same sign at about a quarter of that magnitude, and the
pooled $+0.043$ likewise sits well inside the published interval, so the panel as
a whole is consistent with a positive Shen-Ross uniqueness premium that is real,
heterogeneous across markets, and smaller on average than the single-market
published benchmark would suggest.
